% Part: first-order-logic
% Chapter: models-theories
% Section: size-of-structures

\documentclass[../../../include/open-logic-section]{subfiles}

\begin{document}

\olfileid[tr]{fol}{mat}{siz}

\olsection{\printtoken{P}{structure} Büyüklüğünü İfade Etmek}

\begin{explain}
Bir dilin mantıksal olmayan simgelerini kullanmadan da yapıların bazı
özelliklerini ifade edebiliriz. Örneğin !!a{structure} içinde ancak ve
ancak !!{structure} ifadesinin !!{domain} en az, en çok ya da tam
olarak belirli bir $n$ sayıda !!{element}s içerdiğinde doğru olan
!!{sentence}s vardır.
\end{explain}

\begin{prop}
\begin{multline*}
  !A_{\ge n} \ident \lexists[x_1][\lexists[x_2][\dots\lexists[x_n][{}]]]\\
  \begin{aligned}
  (\eq/[x_1][x_2] \land {}
  \eq/[x_1][x_3] \land \eq/[x_1][x_4] \land \dots \land \eq/[x_1][x_n] \land {}\\
\eq/[x_2][x_3] \land \eq/[x_2][x_4] \land \dots \land {} \eq/[x_2][x_n] \land {} \\
\vdots\\
\eq/[x_{n-1}][x_n])
\end{aligned}
\end{multline*}
!!{sentence} ifadesi, !!a{structure}~$\Struct M$ içinde ancak ve ancak
$\Domain M$ en az $n$ !!{element}s içerdiğinde doğrudur. Dolayısıyla
$\Sat{M}{\lnot !A_{\ge n+1}}$ ancak ve ancak $\Domain M$ en çok $n$
!!{element}s içerdiğinde geçerlidir.
\end{prop}

\begin{prop}
\begin{multline*}
!A_{= n} \ident \lexists[x_1][\lexists[x_2][\dots\lexists[x_n][{}]]] \\
\begin{aligned}
  (\eq/[x_1][x_2] \land {}
  \eq/[x_1][x_3] \land \eq/[x_1][x_4] \land \dots \land \eq/[x_1][x_n] \land {}\\
\eq/[x_2][x_3] \land \eq/[x_2][x_4] \land \dots \land {} \eq/[x_2][x_n] \land {} \\
\vdots\\
\eq/[x_{n-1}][x_n] \land {} \\
\lforall[y][(\eq[y][x_1] \lor \dots \lor \eq[y][x_n]]))
\end{aligned}
\end{multline*}
!!{sentence} ifadesi, !!a{structure}~$\Struct M$ içinde ancak ve ancak
$\Domain M$ tam olarak $n$ !!{element}s içerdiğinde doğrudur.
\end{prop}

\begin{prop}
!!a{structure}, ancak ve ancak
\[
\{!A_{\ge 1}, !A_{\ge 2}, !A_{\ge 3}, \dots \}
\]
kümesinin bir modeliyse sonsuzdur.
\end{prop}

$\Struct M$ içinde ancak ve ancak $\Domain M$ sonsuz olduğunda doğru
olan tek bir salt mantıksal tümce yoktur. Bununla birlikte, yalnızca
sonsuz modelleri olan, mantıksal olmayan !!{predicate}s içeren
!!{sentence}s verilebilir (gerçi her sonsuz !!{structure} bunların bir
modeli değildir). Bir yapının sonlu olması özelliği ve bir yapının
!!{nonenumerable} olması özelliği, sonsuz bir !!{sentence}s kümesiyle
bile ifade edilemez. Bu olgular tıkızlık ve Löwenheim--Skolem
teoremlerinden çıkar.

\end{document}
